\documentclass[professionalfont]{beamer}
\usepackage{xcolor}
\usepackage{adjustbox}
\usepackage{hyperref}
\usepackage{xkeyval}
\usepackage{ragged2e}
\usepackage{amsmath}
\usepackage{amsfonts}
\usepackage{amssymb}
\usepackage{yhmath}
\usepackage{amsthm}
\usepackage{mathrsfs}
\usepackage{mathtools}
\usepackage{mathalfa}
\usepackage{upgreek}
\usepackage{graphics}
\usepackage{yfonts}
\usepackage{dsfont}
\usetheme{Warsaw}
\usecolortheme{wolverine}
%\usetheme{Madrid}
%\usetheme{Bergen}
%\usetheme{AnnArbor}
%\usetheme{CambridgeUS}\setbeamercovered{dynamic}
%\usetheme{Hannover}
%\usetheme{Antibes}
%\usetheme{Copenhagen}
%\usetheme{Marburg}
%\usefonttheme{serif}
\usepackage{xepersian}
\settextfont{XB Niloofar}
\begin{document}
\title{هتل بی‌نهایت هیلبرت}
\author{روح‌الله حسینی‌نوه}
\begin{frame}
\date{}
\maketitle
\begin{center}
بخش ریاضی محض\\
دانشکده ریاضی و کامپیوتر\\
دانشگاه شهید باهنر کرمان

\vspace{10mm}
1403/08/27
\end{center}
\end{frame}
\begin{frame}
\centering
\includegraphics[width=12cm]{2.jpg}
\end{frame}
\begin{frame}
\frametitle{تا چند بلدیم بشمریم؟}
\centering
{\Huge
افراد قبایل هوتانتو تا 3 می‌شمرند اشیاء با تعداد بیشتر از 3 برای آن‌ها

\textbf{شمارش‌ناپذیرند}}
\end{frame}
\begin{frame}

{\Huge جهان از
$10^{82}$
اتم تشکیل شده است.}

\pause

{\huge $$\omega=\lbrace 0, 1, 2, \dots\rbrace$$}

\pause

{\huge $$n \implies n+1$$}

{\huge $$n_i \implies n_{i+1}$$
اثبات حدس گلدباخ، ریتم اعداد اول و ... آسان و ساده نیست
}
\end{frame}
\begin{frame}
\begin{theorem}[اقلیدس]
{\huge اعداد اول نامتناهی هستند!}
بی‌نهایت عدد اول وجود دارد؟
\end{theorem}
\pause
\begin{proof}
فرض کنیم مجموعه متناهی زیر، مجموعه کل اعداد اول است.
$$\lbrace n_1, n_2, \dots, n_k\rbrace$$
$$n_{k+1}=(n_1 \times n_2 \times \dots \times n_k) +1$$
\end{proof}
\pause
\begin{theorem}
{\huge اعداد زوج نامتناهی هستند!}
\end{theorem}
\end{frame}
\begin{frame}
\centering
\includegraphics[height=\paperheight]{1.jpg}
\end{frame}
\begin{frame}

{\huge دیوید هیلبرت در سخنرانی سال 1924 با عنوان
\begin{center}
\lr{Uber das Unendliche}
\end{center}}
\begin{block}{}
آیا اصل لانه کبوتری(دیریکله 1834) همیشه برقرار است؟
\end{block}

\end{frame}
\begin{frame}
\centering
{\Huge
$$A=4.\bar{9}=4.999\dots$$
\pause
$$ 10A=49.\bar{9}$$
\pause
$$10A-A=49.\bar{9}-4.\bar{9}=45$$
\pause
$$4.\bar{9}=5$$
}
\end{frame}
\begin{frame}
\frametitle{ماشین چاپ بی‌نهایت}
{\Large
تمام کتاب‌ها و سخنرانی‌های نیچه، حتی آن‌هایی که دور انداخته است

برعکسشان

یک در میان


کتاب‌هایی که قرار است در آینده چاپ شوند

تمام مقالاتمان
}
\end{frame}
\begin{frame}

{\huge نتایج ضد شهودی به طور اثباتی درست}

\pause
\begin{block}{}
{\LARGE در هتل هیلبرت، عبارات
\begin{center}
"\textbf{همه اتاق‌های هتل پر هستند}"
\end{center}
و
\begin{center}
"\textbf{هیچ مهمان دیگری نمی‌تواند در هتل جا بگیرد}"
\end{center}
معادل نیستند}
\end{block}
\end{frame}
\begin{frame}

{\LARGE هتل بی‌نهایت اتاق دارد و همه اتاق‌ها پر هستند. مهمان جدیدی به هتل وارد می‌شود
\begin{block}{}
$$n \longrightarrow n+1$$
\end{block}
\pause
$m$
مهمان وارد می‌شوند.
\begin{block}{}
$$n \longrightarrow n+m$$
\end{block}
}
\end{frame}
\begin{frame}

{\Large
اتوبوسی که بی‌نهایت صندلی پر دارد به هتل می‌رسد

\begin{block}{}
$$ n \longrightarrow 2n$$
\end{block}

\pause
صف بی‌نهایتی از اتوبوس که هر کدام بی‌نهایت صندلی پر دارند می‌رسد
\begin{block}{}
$$ n \longrightarrow 2^n$$
\end{block}
\begin{block}{}
$$m_{1i} \longrightarrow 3^i$$
$$ m_{2i} \longrightarrow 5^i$$
$$ m_{3i} \longrightarrow 7^i$$
$$\vdots$$
\end{block}
}
\end{frame}
\begin{frame}

{\Huge
در تمام این شب‌ها، ‌درآمد هتل ثابت است!

\pause
\begin{center}
ذره خدا
\end{center}
}
\end{frame}
\begin{frame}

\huge{
گالیله: در جدال فکری خودم با بی‌نهایت به این نتیجه رسیدم که ذهن متناهی انسان قادر به درک نامتناهی نیست.

$$\infty$$
}
\end{frame}
\begin{frame}
\frametitle{اعداد ترامتناهی}

{\Huge



تعداد همه اعداد چقدر است؟

تعداد نقاط یک خط چقدر است؟

\pause

$$\omega= \lbrace 0, 1, 2 , \dots \rbrace $$
\pause

آیا منجر به تناقض می‌شود؟

}
\end{frame}
\begin{frame}
\frametitle{بازگشت به تناظر یک به یک}
{\Huge
چون نمی‌توانیم ترامتناهی را بشمریم باید هوتانتویی رفتار کنیم

\pause
$$f:\mathds{N} \rightarrow \mathds{Z} : \begin{cases}
\frac{n}{2} & \text{زوج}\\
\frac{1-n}{2} & \text{فرد}
\end{cases}
$$

\pause
$$|\mathds{Q}|=|\mathds{N}|= \aleph_0$$
}
\end{frame}
\begin{frame}
\centering
\includegraphics[height=\paperheight]{3.jpg}
\end{frame}
\begin{frame}

{\Huge
جزء با کل برابر است؟ خیر

جزء به بزرگی کل است ( خلاف اصل اقلیدس)

\pause
اصل گسترش ترتیب را بی‌اهمیت می‌کند
$$\lbrace 1, 2, 3 \rbrace = \lbrace 3, 1, 2 \rbrace$$

\pause
کاردینالیته هم همینطور
}
\end{frame}
\begin{frame}
\frametitle{حساب کاردینال‌ها}
{\LARGE
$$n + \aleph_0 = \aleph_0$$
$$\lbrace 0, 1, 2 \dots, n-1 \rbrace \cup \lbrace n, n+1, \dots \rbrace$$
\pause
$$\aleph_0 + \aleph_0 = \aleph_0$$
$$ \lbrace 1, 3, 5 \dots \rbrace \cup \lbrace 0, 2, 4 \dots \rbrace$$
\pause
$$\aleph_0. \aleph_0=\aleph_0$$
}
\end{frame}
\begin{frame}
\frametitle{مجموعه توانی $\omega$}

\Large{
برای هر $A \subseteq \omega$
$$f(n)=\begin{cases}
1 & n \in A\\
0 & n \notin A
\end{cases}
$$
$$A \mapsto 11000010111 \dots$$
\pause
$$2^{\aleph_0}\leq \aleph_0^{\aleph_0}\leq 2^{\aleph_0}$$
}
\end{frame}
\begin{frame}
\includegraphics[height=7cm]{4.jpg}

\lr{ \textit{\textbf{{\large The continuum is not countable}}}}
\end{frame}
\begin{frame}
\begin{theorem}
$$|X| < |\mathcal{P}(X)|$$
\end{theorem}
\begin{corollary}
$$\aleph_0 < 2^{\aleph_0}$$
\end{corollary}
\begin{definition}
$$\aleph_0 < \aleph_1\leq 2^{\aleph_0}$$
\end{definition}
\begin{corollary}
$$\aleph_0, \aleph_1, \aleph_2, \dots, \aleph_{\omega}, \aleph_{\omega+1}, \dots$$
\end{corollary}
\end{frame}
\begin{frame}
\frametitle{به چقدر عدد نیاز داریم؟}
{\huge
$$V_0 = \emptyset$$
$$ V_{n+1}= \mathcal{P}(V_n)$$
$$V_{\omega}= \bigcup_{n \in \omega} V_n$$
}
\begin{block}{}
$$\omega \subset V_{\omega}$$
$$\omega \in V_{\omega+1}$$
$$\mathcal{P}(\omega)=\mathds{R} \in V_{\omega+2}$$
$$\mathcal{P}(\mathcal{P}(\mathds{R})) \in V_{\omega+4}$$
\end{block}
\end{frame}
\begin{frame}
\centering
\includegraphics[width=11cm]{5.jpg}

\lr{ \textit{\textbf{{\large }}}}
\end{frame}
\begin{frame}

\begin{center}
{\Huge \textbf{از توجه شما متشکرم}}
\end{center}
\end{frame}
\end{document}